% ==============================================================================
% ch02_ai_assisted_content.tex — บทที่ 2
% ==============================================================================
\chapter{การใช้ปัญญาประดิษฐ์สกัดตัวชี้วัดหลักสูตรสู่พิมพ์เขียวสื่อการสอน}
\label{chap:ch02}

\begin{teachingplan}{แผนการบริหารการสอนประจำบทที่ 2}
\noindent\textbf{หัวข้อเนื้อหาประจำบท}
\begin{enumerate}[leftmargin=1.5em, itemsep=2pt]
    \item การสังเคราะห์ตัวชี้วัดหลักสูตรแกนกลางสู่ผลลัพธ์การเรียนรู้ระดับพฤติกรรม (Bloom's Revised Taxonomy)
    \item กรอบการสั่งการปัญญาประดิษฐ์ระดับสูงแบบ CLEAR Prompting Framework
    \item เทคนิคการสร้างสถาปัตยกรรมสื่อการสอนเสมือนจริงด้วยโมเดลภาษาขนาดใหญ่ (LLM)
    \item การสร้าง System Prompt และ Role-Based Personas สำหรับครูผู้สร้างสื่อ 3D
    \item การตรวจสอบความถูกต้องทางฟิสิกส์ (Physics Verification) และการป้องกันปัญหา AI Hallucination
\end{enumerate}

\noindent\textbf{วัตถุประสงค์เชิงพฤติกรรม (Behavioral Objectives)}
\begin{enumerate}[leftmargin=1.5em, itemsep=2pt]
    \item ผู้เรียนสามารถแปลงตัวชี้วัดหลักสูตรให้เป็นคำสั่ง Prompt เชิงสถาปัตยกรรมได้อย่างถูกต้อง
    \item ผู้เรียนสามารถใช้กรอบ CLEAR Prompting ในการสร้างโค้ดจำลองฟิสิกส์ HTML/JS ได้
    \item ผู้เรียนสามารถเขียนแบบทดสอบและสคริปต์การตรวจสอบความถูกต้องทางคณิตศาสตร์ร่วมกับ AI ได้
\end{enumerate}

\noindent\textbf{กิจกรรมการเรียนการสอน (Teaching-Learning Activities)}
\begin{enumerate}[leftmargin=1.5em, itemsep=2pt]
    \item อธิบายทฤษฎีการเขียน Prompt เชิงโครงสร้างและการเทียบเคียงกับ Bloom's Taxonomy (30 นาที)
    \item ปฏิบัติการ: ฝึกเขียน CLEAR Prompt เพื่อสั่งการให้ AI ร่างโค้ดจำลองลูกตุ้มเพนดูลัม (45 นาที)
    \item วิเคราะห์ข้อผิดพลาดและวิธีแก้ปัญหาภาพหลอน (Hallucination) ในการคำนวณฟิสิกส์ (45 นาที)
\end{enumerate}

\noindent\textbf{สื่อการเรียนการสอน}
\begin{enumerate}[leftmargin=1.5em, itemsep=2pt]
    \item เอกสารคู่มือบทที่ 2, เทมเพลต CLEAR Prompting, อินเทอร์เฟซผู้ช่วย AI (Claude/Gemini/GPT)
\end{enumerate}

\noindent\textbf{การวัดและประเมินผล}
\begin{enumerate}[leftmargin=1.5em, itemsep=2pt]
    \item ประเมินจากความสมบูรณ์ของชุดคำสั่ง Prompt ที่ผู้เรียนสร้างขึ้น
    \item ประเมินจากความสามารถในการตรวจจับข้อผิดพลาดในโค้ดฟิสิกส์ที่ AI สร้างขึ้น
\end{enumerate}
\end{teachingplan}

\newpage

\section{กรอบการสั่งการปัญญาประดิษฐ์ CLEAR Prompting Framework}
การใช้ Generative AI ในการพัฒนาสื่อการสอนระดับมืออาชีพ ไม่สามารถสั่งการด้วยข้อความสั้นๆ หรือกำกวมได้ ครูและนักพัฒนาจำเป็นต้องใช้ระเบียบวิธีวิศวกรรมคำสั่ง (Prompt Engineering) ที่มีโครงสร้างชัดเจน เรียกว่า **CLEAR Prompting Framework** ประกอบด้วย 5 มิติ:

\begin{figure}[h!]
\centering
\begin{tikzpicture}[scale=0.88, every node/.style={scale=0.88}]
    % 5 Steps of CLEAR
    \node[draw=rbruNavy, fill=blue!10, rounded corners=6pt, text width=2.4cm, align=center, minimum height=2.4cm] (C) at (0,0) {\bfseries\color{rbruNavy}\Large C\\[4pt]\footnotesize Context \& Persona\\(บทบาทและบริบท)};
    \node[draw=rbruBlue, fill=cyan!10, rounded corners=6pt, text width=2.4cm, align=center, minimum height=2.4cm] (L) at (3.2,0) {\bfseries\color{rbruBlue}\Large L\\[4pt]\footnotesize Learning Goal\\(เป้าหมายการเรียนรู้)};
    \node[draw=rbruEmerald, fill=green!10, rounded corners=6pt, text width=2.4cm, align=center, minimum height=2.4cm] (E) at (6.4,0) {\bfseries\color{rbruEmerald}\Large E\\[4pt]\footnotesize Execution Logic\\(ตรรกะเชิงฟิสิกส์)};
    \node[draw=rbruGold, fill=amber!10, rounded corners=6pt, text width=2.4cm, align=center, minimum height=2.4cm] (A) at (9.6,0) {\bfseries\color{rbruGold!80!black}\Large A\\[4pt]\footnotesize Architecture \& UI\\(โครงสร้างส่วนหน้า)};
    \node[draw=rbruPurple, fill=purple!10, rounded corners=6pt, text width=2.4cm, align=center, minimum height=2.4cm] (R) at (12.8,0) {\bfseries\color{rbruPurple}\Large R\\[4pt]\footnotesize Review \& Audit\\(เกณฑ์ตรวจสอบ)};

    \draw[-{Stealth[scale=1.2]}, thick, gray] (C) -- (L);
    \draw[-{Stealth[scale=1.2]}, thick, gray] (L) -- (E);
    \draw[-{Stealth[scale=1.2]}, thick, gray] (E) -- (A);
    \draw[-{Stealth[scale=1.2]}, thick, gray] (A) -- (R);
\end{tikzpicture}
\caption{โครงสร้างกรอบคำสั่ง CLEAR Prompting สำหรับการสร้างสื่อจำลองฟิสิกส์ 3 มิติ}
\label{fig:clear_prompting}
\end{figure}

\begin{enumerate}[leftmargin=1.5em, itemsep=4pt]
    \item \textbf{C -- Context \& Role Persona:} กำหนดบทบาทที่ชัดเจน เช่น \textit{“คุณคือสถาปนิกผู้เชี่ยวชาญการสร้าง 3D Virtual Physics Lab ตามมาตรฐานโอลิมปิกวิชาการ...”}
    \item \textbf{L -- Learning Objectives (Bloom's Taxonomy):} ระบุพฤติกรรมการเรียนรู้ที่ต้องการ เช่น วิเคราะห์การต่อตัวต้านทานแบบผสม (Analyze) และประเมินค่าความต่างศักย์ (Evaluate)
    \item \textbf{E -- Execution \& Physics Formulas:} ใส่สมการคณิตศาสตร์และเงื่อนไขขอบเขตอย่างชัดเจน เช่น $R_{\text{ชุด}} = \frac{r_1 r_2}{r_1 + r_2}$ และต้องทอนเป็นเศษส่วนอย่างต่ำด้วย Euclidean GCD
    \item \textbf{A -- Architecture \& Technology Stack:} ระบุภาษาและไลบรารีที่ต้องการ เช่น Vanilla HTML5, Three.js, MediaPipe, CSS Glassmorphism ปราศจากการใช้ไลบรารีภายนอกที่หนักเกินจำเป็น
    \item \textbf{R -- Robustness \& Constraints:} กฎเหล็กข้อห้าม เช่น ห้ามใช้ค่าสุ่มจำลอง ห้ามทับซ้อนเลย์เอาต์ และต้องได้ผลลัพธ์ 0 Errors / 0 Warnings
\end{enumerate}

\section{เทมเพลตคำสั่ง Master Prompt สำหรับสร้างห้องปฏิบัติการเสมือนจริง}

\begin{aiprompt}{เทมเพลตคำสั่ง CLEAR Prompt สำหรับระบบวงจรไฟฟ้า 3D}
คัดลอกและนำข้อความคำสั่งต่อไปนี้ไปสั่งการในโมเดลปัญญาประดิษฐ์:
\begin{lstlisting}
[ROLE & CONTEXT]:
ทำหน้าที่เป็น Senior 3D Physics Simulation Architect พัฒนาระบบห้องปฏิบัติการเสมือนจริงสำหรับนักเรียนโอลิมปิกฟิสิกส์ มหาวิทยาลัยราชภัฏรำไพพรรณี

[LEARNING OUTCOME]:
นักเรียนสามารถเข้าใจการลดรูปวงจรไฟฟ้าตัวต้านทานรูปหกเหลี่ยมด้านเท่าที่มีขั้วต่อ a, b, c, d, e, f และสังเกตการไหลของกระแสไฟฟ้ารวม I และแรงดันคร่อมโพรบ V_PQ

[PHYSICS & MATH CONSTRAINTS]:
1. ตัวต้านทาน 6 ตัว: 1, 2, 3, 4, 5, 6 โอห์ม เรียงตามลำดับ ผลรวม 21 โอห์ม
2. แหล่งจ่ายไฟ DC 10V และตัวต้านทานอนุกรมคงที่ 10 โอห์ม
3. เมื่อเลือกคู่โพรบ P, Q วงแหวนต้องคลี่เป็น 2 กิ่งขนาน r1 และ r2 ทันที
4. คำนวณ R_ชุด = (r1 * r2) / (r1 + r2) ด้วย Euclidean GCD เป็นเศษส่วนแท้
5. คำนวณ I = 10 / (10 + R_ชุด) และ V_PQ = I * R_ชุด

[TECH STACK & ARCHITECTURE]:
- WebGL 3D ด้วย Three.js PBR Shaders
- เอนจินวาดผังลดรูปวงจร 4 ขั้นตอนด้วยเวกเตอร์ SVG พลวัต
- แผงควบคุม Glassmorphism ซ้าย-ขวา ย่อขยายและพับเก็บได้
- โค้ดทั้งหมดต้องเป็น Standalone รันบนเบราว์เซอร์ได้ทันทีโดยไม่ต้อง Build
\end{lstlisting}
\end{aiprompt}

\section{การตรวจสอบความถูกต้องและป้องกันปัญหาภาพหลอน (AI Hallucination)}
ปัญหาที่พบบ่อยเมื่อใช้ AI สร้างโค้ดวิทยาศาสตร์ คือ AI มักประมาณค่าทศนิยม หรือสร้างฟังก์ชันจำลองที่มีลักษณะสุ่ม (Random) เพื่อให้เข็มมิเตอร์ขยับได้ โดยไม่ได้คำนวณจากกฎของฟิสิกส์จริง

\section{สถาปัตยกรรมเอเจนต์ปัญญาประดิษฐ์ผู้ช่วยอัจฉริยะ (Mr. Deep\_xAI \& J.A.R.V.I.S. Protocol)}
ในระบบ RBRU Next Gen Lab 2026 ผู้เรียนไม่ได้ทำการทดลองเพียงลำพัง แต่มี **Mr. Deep\_xAI** ซึ่งเป็น Agentic AI Co-Creation Bot ประจำห้องแล็บ ทำหน้าที่เป็นผู้ช่วยอัจฉริยะแบบเรียลไทม์ โดยมีคุณลักษณะเด่นเชิงสถาปัตยกรรมดังนี้:

\begin{figure}[h!]
\centering
\begin{tikzpicture}[scale=0.85, every node/.style={scale=0.85}]
    % Agentic AI Core Architecture Diagram
    \node[draw=rbruNavy, fill=blue!15, rounded corners=8pt, text width=3.4cm, align=center, minimum height=2.6cm] (VISION) at (0,3) {
        \bfseries\color{rbruNavy} 📷 Multi-Modal Vision\\
        \footnotesize WebGL Viewport Snapshot\\
        ตรวจจับอุปกรณ์ในฉาก 3 มิติ\\
        พิกัด 3 มิติ \& สถานะสายไฟ
    };

    \node[draw=rbruEmerald, fill=green!15, rounded corners=8pt, text width=3.4cm, align=center, minimum height=2.6cm] (BRAIN) at (5.5,3) {
        \bfseries\color{rbruEmerald} 🧠 Agentic Co-Creation\\
        \footnotesize Gemini 2.5 Flash LLM\\
        วิเคราะห์ผลการทดลอง\\
        แนะนำสมการและตรวจแบบทดสอบ
    };

    \node[draw=rbruGold, fill=amber!15, rounded corners=8pt, text width=3.4cm, align=center, minimum height=2.6cm] (VOICE) at (11,3) {
        \bfseries\color{rbruGold!80!black} 🎙️ Neural Male TTS\\
        \footnotesize Formant Pitch (0.44 / 0.58)\\
        J.A.R.V.I.S. Iron Man Protocol\\
        เสียงหนุ่มสุโขทัยธรรมชาติ
    };

    \node[draw=cyan!80!black, fill=cyan!15, rounded corners=8pt, text width=8.5cm, align=center, minimum height=1.6cm] (HUD) at (5.5,-0.2) {
        \bfseries\color{cyan!80!black} ⚡ Web Audio API Iron Man Communicator DSP Filter\\
        \footnotesize High-Pass (120 Hz) + Peaking Presence (3400 Hz, +2.6 dB) + Dynamics Compressor + HUD Blip Chime
    };

    \draw[-{Stealth[scale=1.2]}, line width=1.5pt, gray] (VISION) -- (BRAIN);
    \draw[-{Stealth[scale=1.2]}, line width=1.5pt, gray] (BRAIN) -- (VOICE);
    \draw[-{Stealth[scale=1.2]}, line width=1.5pt, cyan] (VOICE) -- (HUD);
\end{tikzpicture}
\caption{สถาปัตยกรรมระบบปัญญาประดิษฐ์เอเจนต์ Mr. Deep\_xAI (J.A.R.V.I.S. Protocol)}
\label{fig:ai_agent_arch}
\end{figure}

\subsection{การมองเห็นฉาก 3 มิติ (3D Scene Vision Audit)}
Mr. Deep\_xAI สามารถ "มองเห็น" สิ่งที่เกิดขึ้นบนหน้าจอการทดลองของผู้เรียนได้ทันที โดยเมื่อผู้เรียนกดปุ่ม \textbf{📷 สแกน 3D} ระบบจะดึงเฟรมเรนเดอร์สดจาก WebGL Viewport แปลงเป็นภาพ Base64 JPEG ส่งเข้าสู่ Multimodal Vision Engine เพื่อตรวจสอบ:
\begin{enumerate}[leftmargin=1.5em, itemsep=2pt]
    \item ตำแหน่งของโพรบวัด $P$ และ $Q$ ว่าต่ออยู่ที่ขั้วใดของวงจร 6 เหลี่ยม
    \item สภาวะของสวิตช์เปิด-ปิด และค่าแรงดันที่ตั้งไว้บนแหล่งจ่ายไฟ DC
    \item ความยาวเชือกลูกตุ้ม $L$ และการแกว่งมุม $\theta$ ในปฏิบัติการเพนดูลัม
\end{enumerate}

\subsection{ระบบเสียงสังเคราะห์ชายแท้และฟิลเตอร์ชุดเกราะไอรอนแมน (Iron Man Communicator DSP)}
เพื่อให้เสียงบรรยายมีบุคลิกน่าเชื่อถือ สุภาพ และเหมือนผู้ช่วย AI ระดับไฮเอนด์ ระบบได้ออกแบบกลไกเสียง 2 รูปแบบ:
\begin{itemize}[leftmargin=1.5em, itemsep=3pt]
    \item \textbf{โปรโตคอล J.A.R.V.I.S. (เสียงผู้ชายทุ้มลึก):} ปรับจูนระดับเสียงหลักลงสู่ย่านเสียงหน้าอก (\textbf{Pitch: $0.44$}, \textbf{Rate: $0.90$}) เพื่อให้ได้น้ำเสียงทุ้ม สุขุม หนักแน่น เหมือน Paul Bettany ในภาพยนตร์ Iron Man พร้อมสำเนียงการตอบรับแบบสุภาพบุรุษ (เช่น ``ครับท่าน (Sir)'', ``ยินดีรับใช้ครับท่าน'')
    \item \textbf{โปรโตคอลหนุ่มสุโขทัย (เสียงผู้ชายวัยรุ่นเป็นกันเอง):} ปรับระดับเสียง (\textbf{Pitch: $0.58$}, \textbf{Rate: $0.96$}) ผสานคำลงท้ายสำเนียงสุโขทัย (เช่น ``เน้อครับเน้อ'')
    \item \textbf{วงจร DSP จำลองสัญญาณสื่อสารในชุดเกราะ (Web Audio API):} สัญญาณเสียงสังเคราะห์จะถูกส่งผ่านตัวกรองย่านความถี่ต่ำ \textbf{High-Pass Filter ($120\text{ Hz}$)}, ตัวเสริมความคมชัดประกายเสียง \textbf{Peaking Presence Filter ($3400\text{ Hz}$, Gain $+2.6\text{ dB}$)} และ \textbf{Dynamics Compressor} พร้อมเสียงสังเคราะห์แจ้งเตือน \textbf{HUD Blip Chime ($880\text{ Hz} \rightarrow 1320\text{ Hz}$)} ก่อนการพูดทุกครั้ง
\end{itemize}

\begin{tipbox}{เทคนิคการตรวจสอบความถูกต้องทางฟิสิกส์ (Physics Sanity Check)}
\begin{itemize}[leftmargin=1.5em, itemsep=2pt]
    \item \textbf{ตรวจสอบกฎการอนุรักษ์พลังงานและประจุ (KCL \& KVL):} กระแสที่ไหลเข้าจุดต่อ $P$ ต้องเท่ากับผลรวมกระแสในกิ่งที่ 1 และกิ่งที่ 2 เสมอ ($I_{\text{รวม}} = I_1 + I_2$)
    \item \textbf{ตรวจสอบกรณีสุดโต่ง (Boundary Conditions):} เมื่อสวิตช์เปิดอยู่ วงจรต้องมีค่า $I = 0\text{ A}$ และ $V_{PQ} = 0\text{ V}$ เสมอ
    \item \textbf{ตรวจสอบการทอนเศษส่วน:} ตรวจสอบว่า $\frac{20}{21}$ โอห์ม เมื่อต่ออนุกรมกับ $10$ โอห์ม จะได้ $\frac{230}{21}$ โอห์ม ตรงตามทฤษฎีคณิตศาสตร์แท้จริงหรือไม่
\end{itemize}
\end{tipbox}

\section*{คำถามทบทวนและใบงานปฏิบัติการท้ายบทที่ 2}
\addcontentsline{toc}{section}{คำถามทบทวนและใบงานปฏิบัติการท้ายบทที่ 2}

\begin{enumerate}[leftmargin=1.5em, itemsep=8pt]
    \item จงอธิบายองค์ประกอบของกรอบ CLEAR Prompting และชี้แจงว่าเพราะเหตุใดขั้นตอน \textbf{E (Execution Logic)} จึงมีความสำคัญอย่างยิ่งยวดสำหรับสื่อการสอนวิทยาศาสตร์
    \item หาก AI สร้างโค้ดคำนวณความต้านทานสมมูลแบบขนานมาเป็น $R = \frac{r_1 + r_2}{2}$ ท่านจะเขียนคำสั่ง Prompt อย่างไรเพื่อตักเตือนและชี้นำให้ AI แก้ไขเป็นสมการที่ถูกต้อง
    \item \textbf{ใบงานปฏิบัติการ 2.1:} ให้นักศึกษาใช้กรอบ CLEAR Prompting ร่างคำสั่งฉบับสมบูรณ์สำหรับสร้างการทดลอง “การแกว่งของลูกตุ้มเพนดูลัมฮาร์มอนิกอย่างง่ายเพื่อหาค่าความเร่ง $g$” โดยต้องระบุสมการ $T = 2\pi\sqrt{L/g}$ และเงื่อนไขการวาดกราฟ $T^2$ กับ $L$ ด้วยวิธีกำลังสองน้อยที่สุด (OLS)
\end{enumerate}
